\documentstyle[%


righttag%


,11pt%


,amssymb%


,russian%               remove in Eng.


,izvru%                 Set izven% in Eng.


% ,izvru-ks             for brief communications


]{amsart}





%       Math definitions





\newcommand{\la}{\langle}


\newcommand{\bsl}{\boldsymbol [}


\newcommand{\bsr}{\boldsymbol ]}


\newcommand{\ra}{\rangle}


\newcommand{\ges}{\geqslant}


\newcommand{\les}{\leqslant}


\newcommand{\Con}{\operatorname{Con}}


\newcommand{\St}{\operatorname{St}}


\newcommand{\tts}[1]{}








\theoremstyle{definition}


\theorembodyfont{\rm}


\newtheorem{cornonum}{\hskip\parindent Следствие}


\newtheorem{defnnonum}{\hskip\parindent Определение}


\renewcommand{\thecornonum}{}


\renewcommand{\thedefnnonum}{}





%\hoffset=-15mm                  For Eng.





\setcounter{page}{52}


\god{1997}


\nomer{11~(426)} %                Remove in Eng.


%\nomer{Vol.\,41, No.\,11}        Set in Eng.


% \str{75--81}                    Set in Eng.


\udk{519.48}


\author{\it А.Г.\,Пинус}


\title{Об $\aleph$-идентичности отношений эпиморфности\\


и вложимости на квазимногообразиях}





\thanks{Работа выполнена при поддержке Российского фонда фундаментальных


исследований (код проекта 94-01-00183).}


\date{}


% \pagestyle{myheadings}         Set in Eng.





\pagestyle{plain} %              Remove in Eng.


\begin{document}





% \thispagestyle{plain}          Set in Eng.





\maketitle





% Body text








Работа посвящена вопросу взаимообусловленности отношений эпиморфности и


вложимости на алгебрах некоторых многообразий (квазимногообразий). В


работе автора \cite{P1} доказана финитарная независимость отношений


эпиморфности и вложимости на алгебрах любого нетривиального


конгруэнц-дистрибутивного многообразия с продолжимыми конгруэнциями. С


другой стороны, представляет интерес вопрос существования больших


совокупностей алгебр рассматриваемого класса, на которых отношения


эпиморфности и вложимости совпадают. В \cite{P2} показано, что для любого


кардинала $\aleph\ges\aleph_0$ в любом нетривиальном


конгруэнц-дистрибутивном многообразии


$\frak M$ с продолжимыми конгруэнциями найдется совокупность


$\frak M$-алгебр,


имеющих


мощность $2^{\aleph}$, такая, что отношения эпиморфности и вложимости


между


алгебрами этой совокупности совпадают и при этом сама совокупность с


этими отношениями изоморфна множеству всех подмножеств некоторого


множества мощности $2^{\aleph}$ с отношением теоретико-множественного


включения.


Известно, что для счетных $\frak M$-алгебр подобный результат, вообще


говоря,


не верен. В данной работе приведены некоторые достаточные условия, при


которых аналогичный результат имеет место для счетных алгебр ряда


многообразий (квазимногообразий).





Напомним, что если $\frak K$ --- некоторый класс универсальных алгебр,


то через


$\frak I\frak K$ обозначается совокупность типов изоморфизма


$\frak K$-алгебр,


если $\aleph$ ---


произвольный кардинал, то


$\frak K_{\aleph}=\{\frak A\in\frak K\bigr|\bigl.|\frak A|\les\aleph\}$.


На совокупности $\frak I\frak K$ определим два отношения


квазипорядка $\les$ (вложимости) и $\ll$ (эпиморфности) следующим


образом: для


$a,b\in\frak I\frak K$ отношение $a\les b$ ($a\ll b$) имеет место тогда


и только тогда, когда алгебра типа


$a$ изоморфно вложима в алгебру типа $b$ (является гомоморфным образом


алгебры типа $b$). {\it Скелетом}


эпиморфности (вложимости) класса $\frak K$ называется


квазиупорядоченный класс $\la\frak I\frak K;\ll\ra$


($\la\frak I\frak K;\les\ra$), {\it счетным скелетом}


эпиморфности


(вложимости) класса $\frak K$ --- квазиупорядоченное множество


$\la\frak I\frak K_{\aleph_0};\ll\ra$


($\la\frak I\frak K_{\aleph_0};\les\ra$). {\it Двойным


скелетом} ({\it двойным счетным скелетом}) класса $\frak K$ является


дважды квазиупорядоченный класс (множество)


$\la\frak I\frak K;\ll,\les\ra$


($\la\frak I\frak K_{\aleph_0};\ll,\les\ra$). Более подробно об этих


понятиях и о результатах, с ними связанных, см. \cite{P2}, \cite{P3},


\cite{P4}. Далее рассматриваем лишь алгебры не более чем счетной


сигнатуры.





Отношения эпиморфности и вложимости на классе $\frak K$ универсальных


алгебр


называются {\it финитарно} ({\it локально} \cite{P1})


{\it независимыми}, если


любое


конечное дважды квазиупорядоченное множество изоморфно вложимо в двойной


скелет $\la\frak I\frak K;\ll\les\ra$ класса $\frak K$. В \cite{P1}


доказана финитарная независимость


отношений эпиморфности и вложимости на любом нетривиальном


конгруэнц-дистрибутивном многообразии с продолжимыми конгруэнциями. Более


точно, там же доказано, что для любого такого многообразия $\frak M$


произвольное конечное дважды квазиупорядоченное множество


$\la A;\les_1,\les_2\ra$ изоморфно


вложимо в $\la\frak I\frak M_{\aleph_1};\ll,\les\ra$. Согласно же


результатам \cite{P5} нельзя утверждать


подобную вложимость любого конечного дважды квазиупорядоченного множества


$\la A;\les_1,\les_2\ra$ в счетный скелет


$\la\frak I\frak M_{\aleph_0};\ll\les\ra$ произвольного нетривиального


конгруэнц-дистрибутивного многообразия с продолжимыми конгруэнциями (даже


произвольного дискриминаторного многообразия), т.е. отношения


эпиморфности и вложимости на счетных алгебрах подобных многообразий уже


не обязаны быть финитарно независимыми.





Пусть $\la A;\les\ra$ --- некоторое квазиупорядоченное множество, через


$\la A;\les,\les\ra$


обозначим


множество $A$, наделенное двумя квазипорядками, второй из которых


идентичен первому.





\begin{defnnonum}


Отношения эпиморфности и вложимости на классе $\frak K$ универсальных


алгебр


назовем $\la A;\les\ra$-{\it идентичными}, если существует изоморфное


вложение дважды


квазиупорядоченного множества $\la A;\les,\les,\ra$ в двойной скелет


$\la\frak I\frak K;\ll,\les\ra$


класса $\frak K$. Отношения


эпиморфности и вложимости на классе $\frak K$ назовем


$\aleph$-{\it идентичными}


({\it частично $\aleph$-идентичными}),


если эти отношения $\la A,\les\ra$-идентичны на классе


$\frak K_{\aleph}$


для любого


квазиупорядоченного (частично упорядоченного) множества


$\la A;\les\ra$ мощности $\aleph$.


\end{defnnonum}





Очевидно, что если $\la P(\aleph);\subseteq\ra$ --- совокупность всех


подмножеств кардинала $\aleph$,


частично упорядоченная отношением теоретико-множественного включения, то


из $\la P(\aleph);\subseteq\ra$-идентичности отношений вложимости и


эпиморфности на


классе $\frak K_{\aleph}$ следует частичная $\aleph$-идентичность этих


отношений на классе $\frak K$.





Как отмечено в начале работы, для любого нетривиального


конгруэнц-дистрибутивного многообразия $\frak M$ с продолжимыми


конгруэнциями,


для любого $\aleph\ges\aleph_0$ отношения эпиморфности и вложимости на


$\frak M$


частично


$2^{\aleph}$-идентичны. Ниже будут приведены некоторые достаточные


условия


$\aleph_0$-идентичности отношений эпиморфности и вложимости на


квазимногообразиях (многообразиях).





В дальнейшем через $\Con\frak A$ будем обозначать решетку конгруэнций


алгебры


$\frak A$. Для любого класса $\frak K$, для любой алгебры $\frak A$


конгруэнцию $\theta\in\Con\frak A$ назовем


$\frak K$-конгруэнцией, если $\frak A/\theta\in\frak K$. Совокупность


$\frak K$-конгруэнций на алгебре $\frak A$,


частично упорядоченную отношением включения, будем обозначать через


$\Con^{\frak K}\frak A$.





Для любого класса алгебр $\frak K$ неодноэлементная $\frak K$-алгебра


$\frak A$


называется


$\frak K$-простой, если для любой нетривиальной конгруэнции $\theta$


(т.е. отличной


от наибольшей $\bigtriangledown_{\frak A}$ и наименьшей


$\bigtriangleup_{\frak A}$) алгебры $\frak A$


фактор-алгебра $\frak A/\theta$


не принадлежит $\frak K$. Если $\frak K$ замкнут относительно


гомоморфных образов, то


очевидно, что понятие $\frak K$-простой алгебры совпадает с понятием


простой


алгебры. Напомним, что по теореме Магари любое нетривиальное многообразие


$\frak M$ содержит не более чем счетную алгебру. Это утверждение


В.А.Горбунов


\cite{Gor} обобщил для квазимногообразий: любое нетривиальное


квазимногообразие $\frak K$ содержит не более чем счетную


$\frak K$-простую


алгебру.





Класс $\frak K$ обладает свойством Фрезера-Хорна, если для любых


$\frak K$-алгебр


$\frak A_1$,


$\frak A_2$ и любой $\frak K$-конгруэнции $\theta$ на алгебре


$\frak A_1\times\frak A_2$ найдутся


$\frak K$-конгруэнции $\theta_i\in\Con^{\frak K}\frak A_i$ такие, что


$\theta=\theta_1\times\theta_2$.





Имеет место следующий довольно очевидный факт.





\begin{thm}


Если $\frak K$ --- квазимногообразие, обладающее свойством Фрезера-Хорна


и $\frak K$ содержит $\aleph$ попарно невложимых друг в друга


$\frak K$-простых


конечно-порожденных алгебр, обладающих одноэлементными подалгебрами, то


отношения эпиморфности и вложимости на


$\frak K\;$ $\aleph$-идентичны.


\end{thm}





\begin{pf}


Пусть $\frak A_i$ ($i\in\aleph$) --- попарно невложимые друг в друга


$\frak K$-простые


конечно-порожденные алгебры и $\{e_i\}$ --- одноэлементные подалгебры


алгебр


$\frak A_i$. Для любого $B\subseteq\aleph\setminus\{0\}$ пусть


$\frak A_B$ ---


прямая сумма алгебр $\frak A_i$ ($i\in B$), т.е.


$\frak A_B=\{f\in\prod\limits_{i\in B}\frak A_i\bigl|\bigr.


|\{j\in B\mid f(j)\ne e_j\}|<\aleph_0\}$.





Для доказательства частичной $\aleph$-идентичности достаточно заметить,


что для


любых $B_1,B_1\subseteq\aleph$ отношения


$\frak A_{B_1}\les\frak A_{B_2}$ ($\frak A_{B_1}\ll\frak A_{B_2}$) имеют


место тогда и только тогда, когда $B_1\subseteq B_2$.


Если $B_1\subseteq B_2$, то очевидно, что алгебра $\frak A_{B_1}$


изоморфно


вложима в алгебру $\frak A_{B_2}$ и является гомоморфным образом алгебры


$\frak A_{B_2}$.





Докажем обратное. Пусть теперь $\varphi$ --- изоморфное вложение


некоторой


алгебры $\frak A_{B_1}$ в алгебру $\frak A_{B_2}$ и $j\in B_1$. Так как


алгебра $\frak A_j$ изоморфно вложима в


алгебру $\frak A_{B_1}$, то существует вложение алгебры $\frak A_j$ в


алгебру $\frak A_{B_2}$. В силу же


простоты алгебры $\frak A_j$ найдется $e\in B_2$ такое, что


$\frak A_j$ вложима в


алгебру $\frak A_e$. По


условию же алгебры $\frak A_i$ ($i\in\aleph$) попарно не вложимы друг в


друга. Отсюда $j=e$,


т.е. $B_1\subseteq B_2$ в случае вложимости алгебры


$\frak A_{B_1}$ в алгебру $\frak A_{B_2}$. Если


теперь


$\psi$ --- гомоморфизм алгебры $\frak A_{B_1}$ на алгебру


$\frak A_{B_2}$ и $j\in B_1$, то


ввиду конечной


порожденности алгебры $\frak A_j$ найдутся конечное подмножество $C$


множества


$B_2$ и некоторый гомоморфизм $\psi'$ алгебры


$\prod\limits_{e\in C}\frak A_e$ на алгебру $\frak A_j$. В силу


свойства


Фрезера-Хорна для квазимногообразия $\frak K$ найдутся гомоморфизмы


$\psi'_e$


алгебр $\frak A_e$


на некоторые $\frak K$-алгебры $\frak A'_e$ такие, что


$\psi'=\prod\limits_{e\in C}\psi'_e$ и


$\frak A_j=\prod\limits_{e\in C}\frak A'_e$. Ввиду


$\frak K$-простоты алгебры $\frak A_j$


алгебра $\frak A_j$ изоморфна одной из алгебр


$\frak A'_e$ ($e\in C$), к примеру, $\frak A'_{e_0}$ ($e_0\in C$).


Но $\frak K$-простота алгебры $\frak A_{e_0}$ влечет изоморфность


$\frak A_j$ и $\frak A_{e_0}$.


Попарная же


невложимость друг в друга алгебр $\frak A_i$ ($i\in\aleph$) приводит к


тому,


что для любого $j\in B_1$ имеет место включение $j\in B_2$, т.е.


$B_1\subseteq B_2$ в


случае, если $\frak A_{B_1}\ll\frak A_{B_2}$.





Для доказательства $\aleph$-идентичности отношений вложимости и


эпиморфности на


$\frak K$ достаточно теперь ``размыть'' вложения дважды частично


упорядоченных


множеств $\la A;\les,\les\ra$ в двойной скелет


$\la\frak I\frak K_{\aleph};\les,\ll\ra$ до вложений дважды


квазиупорядоченных


множеств $\la A;\les,\les\ra$ в скелет


$\la\frak I\frak K_{\aleph};\les,\ll\ra$. Для этого достаточно построить


некоторую


совокупность $\frak K_{\aleph}$-алгебр $D$ такую, что


$|D|=\aleph$, алгебры из $D$


попарно


неизоморфны друг другу, для различных


$\frak C_1,\frak C_2\in D\;$ $\frak C_1\ll\frak C_2$,


$\frak C_1\les\frak C_2$ и такую, что для


любых $B_1,B_2\subseteq\aleph$


алгебры $\frak A_{B_1}\times\frak C_1$,


$\frak A_{B_2}\times\frak C_2$ неизоморфны, а отношения


$\frak A_{B_1}\times\frak C_1\ll\frak A_{B_2}\times\frak C_2$


($\frak A_{B_1}\times\frak C_1\les\frak A_{B_2}\times\frak C_2$) для


построенных выше


алгебр $\frak A_B$ имеют место тогда и только тогда, когда


$\frak A_{B_1}\ll\frak A_{B_2}$


($\frak A_{B_1}\les\frak A_{B_2}$).





Хорошо известно (см., напр, \cite{P1}) существование для любого


$\aleph\ges\aleph_0$


совокупности $D'$ попарно неизоморфных булевых алгебр мощности $\aleph$


таких,


что $|D'|=\aleph$ и для $\frak B_1,\frak B_2\in D'$


$\frak B_1\ll\frak B_2$, $\frak B_1\les\frak B_2$. Нетрудно заметить,


что булевы степени


$\frak A^{\frak B_1}_0$, $\frak A^{\frak B_2}_0\;$ $\frak K$-простой алгебры


$\frak A_0$ для неизоморфных булевых алгебр $\frak B_1$, $\frak B_2$ и


сами остаются неизоморфными. Очевидно также, что отношения


$\frak B_1\ll\frak B_2$ ($\frak B_1\les\frak B_2$) влекут


отношения $\frak A_0^{\frak B_1}\ll\frak A_0^{\frak B_2}$


($\frak A_0^{\frak B_1}\les\frak A_0^{\frak B_2}$).





С другой стороны, рассуждения, аналогичные приведенным в начале


доказательства теоремы, показывают, что отношения


$\frak A_{B_1}\times\frak A_0^{\frak B_1}\ll\frak A_{B_2}\times


\frak A_0^{\frak B_2}$


($\frak A_{B_1}\times\frak A_0^{\frak B_1}\les\frak A_{B_2}\times


\frak A_0^{\frak B_2}$)


имеют место


тогда и только тогда, когда $\frak A_{B_1}\ll\frak A_{B_2}$


($\frak A_{B_1}\les\frak A_{B_2}$) и алгебры


$\frak A_{B_1}\times\frak A_0^{\frak B_1}$,


$\frak A_{B_2}\times\frak A_0^{\frak B_2}$ для различных $\frak B_1$ и


$\frak B_2$ неизоморфны. Поэтому в качестве требуемой совокупности $D\;$


$\frak K$-алгебр


годится совокупность $\{\frak A_0^{\frak B}|\frak B\in D'\}$, и


$\aleph$-идентичность


отношений вложимости и эпиморфности на квазимногообразии


$\frak K$ доказана.


\end{pf}





Автор \cite{P7} ввел понятие бедной алгебры. Через


$\frak D_2$ обозначим обычное


двоичное дерево, т.е. совокупность всех конечных кортежей из нулей и


единиц, частично упорядоченную отношением $\les$, где


$\Bar a\les\Bar b$, если кортеж $\Bar b$


является начальным интервалом кортежа $\Bar a$. Под системой уравнений


над


алгеброй $\frak A$ будем понимать конечную совокупность $\frak T$


равенств


$$


\frak T=


\begin{cases}


t^1_1(x,\Bar a)=t^2_1(x,\Bar a),\\


\dots \\


t^1_n(x,\Bar a)=t^2_n(x,\Bar a),


\end{cases}


$$


где $t^i_j(x,\Bar y)$ --- некоторые термы алгебры $\frak A$, а


$\Bar a$ ---


кортеж элементов


этой алгебры. {\it Совокупностью решений} $S_{\frak A}(\frak T)$ системы


уравнений $\frak T$ в алгебре $\frak A$


назовем множество


$\{b\in\frak A\mid\frak A\vDash\overset{n}{\underset{i=1}{\&}}{}


t^1_i(b,\Bar a)=t^2_i(b,\Bar a)\}$.


Очевидно, что множество $S^{\frak A}$ совокупностей решений


всех систем уравнений над алгеброй является нижней полурешеткой


$\frak S_{\frak A}=\la S^{\frak A};\cap\ra$


относительно теоретико-множественного пересечения. Согласно \cite{P7}


алгебра $\frak A$ называется {\it бедной}, если не существует


дизъюнктного вложения


частичного порядка $\frak D_2$ в частичный порядок


$\la S^{\frak A};\subseteq\ra$. При этом под


дизъюнктным


вложением $\frak D_2$ в $\la S^{\frak A};\subseteq\ra$ понимаем такое


изоморфное вложение частично упорядоченного


множества $\frak D_2$ в $\la S^{\frak A};\subseteq\ra$, что для любых


несравнимых $a,b\in\frak D_2$ подмножества


$\varphi(a)$, $\varphi(b)$ основного множества алгебры $\frak A$


дизъюнктны, т.е. $\varphi(a)\cap\varphi(b)=\varnothing$. Алгебру, не


являющуюся бедной, назовем {\it богатой}.





\begin{thm}


Если $\frak K$ --- квазимногообразие, содержащее конечно-порожденную


богатую $\frak K$-простую алгебру, и класс $\frak K$-простых алгебр


универсально


аксиоматизируем, то существует $2^{\aleph_0}$ попарно невложимых друг в


друга конечно-порожденных $\frak K$-простых алгебр.


\end{thm}





\begin{pf}


В дальнейшем для любых алгебр $\frak A_i$ ($i\in I$), для


$f,g\in\prod\limits_{i\in I}\frak A_i$ через


$\bsl f=g\bsr$


обозначим множество $\{i\in I\mid f(i)=g(i)\}$, а через


$\bsl f\ne g\bsr$ --- его


дополнение $I\setminus\bsl f=g\bsr$.





Пусть $\frak A$ --- конечно-порожденная богатая $\frak K$-простая


алгебра. Перечислим


без повторений элементы $\frak A$:


$a_0,a_1,\dotsc,a_n,\dotsc$ Для простоты будем считать, что $\frak A$


порождена двумя элементами $a_0$ и $a_1$. Пусть $\frak A'$


--- подалгебра алгебры $\frak A^{\omega}$,


порожденная элементами $f_0$, $f_1$ и $g$, где для любых


$i\in\omega$ и $j\in\omega\;$ $f_i(j)=a_i$ и $g(j)=a_j$.


Если $I\subseteq\omega$, то через $\pi_I$ обозначим естественное


проектирование алгебры $\frak A'$ на


подалгебру алгебры $\frak A^I$. Заметим, что любое уравнение на


$\frak A$ эквивалентно


(т.е. имеет те же решения) уравнению вида


$t^1(x,a_0,a_1)=t^2(x,a_0,a_1)$, где $t^1$, $t^2$ ---


некоторые термы сигнатуры квазимногообразия $\frak K$ от переменных


$x$, $y_0$, $y_1$.


Заметим также, что любые элементы алгебры $\frak A'$ представимы в виде


$t(g,f_0,f_1)$, где


$t$ --- трехместный терм. Отождествляя натуральное число $j\in\omega$ с


элементом $a_j\in\frak A$


для любых термов $t^1(x,y_0,y_1)$, $t^2(x,y_0,y_1)$ сигнатуры алгебры


$\frak A$,


получим очевидное равенство


$$


\bsl t^1(g,f_0,f_1)=t^2(g,f_0,f_1)\bsr =


S_{\frak A}(t^1(x,a_0,a_1)=t^2(x,a_0,a_1)).


$$


Точно так же для любой системы уравнений над $\frak A$


$$


\frak T=


\begin{cases}


t^1_1(x,a_0,a_1)=t^2_1(x,a_0,a_1),\\


\dots \\


t^1_n(x,a_0,a_1)=t^2_n(x,a_0,a_1)


\end{cases}


$$


для любого $I\subseteq\omega$ выполняется


$$


\pi_I(\frak A')\vDash\overset{n}{\underset{i=1}{\&}}{}


t^1_i(g,f_0,f_1)=t^2_i(g,f_0,f_1)


$$


тогда и только тогда, когда $I\subseteq S_{\frak A}(\frak T)$.





Под фильтром $\frak F$ на нижней полурешетке


$\la L;\wedge\ra$ понимаем такую


совокупность


элементов из $L$, что если $a\in\frak F$, $b\ges a$, то


$b\in\frak F$, и для любых $a,c\in\frak F$


элемент $a\wedge c$ также


принадлежит $\frak F$. Ультрафильтр на $\la L;\wedge\ra$ --- это


максимальный по включению


среди фильтров на $L$, отличных от самой полурешетки $L$. Очевидно, что


любой ультрафильтр $\frak F$ на полурешетке $\frak S_{\frak A}$


совокупностей решений систем


уравнений над алгеброй $\frak A$ (совокупность ультрафильтров на


$\frak S_{\frak A}$


обозначим


через $\frak S^*_{\frak A}$) содержится в некотором ультрафильтре


$\frak F'$ над


$\omega$. Точнее, если $\St(\omega)$


--- совокупность всех ультрафильтров на $\omega$, т.е. стоуновское


пространство


булевой алгебры всех подмножеств множества $\omega$, то через


$\Check{\frak F}$ (для


ультрафильтра $\frak F$ на полурешетке $\frak S_{\frak A}$) обозначим


$\{\frak G\in\St(\omega)\mid\frak G\supseteq\frak F\}$.


Очевидно, что $\Check{\frak F}$ ---


замкнутое подмножество пространства $\St(\omega)$ и для различных


ультрафильтров


$\frak F_1$, $\frak F_2$ на полурешетке $\frak S_{\frak A}$ множества


$\Check{\frak F}_1$, $\Check{\frak F}_2$ дизъюнктны. Для каждого


ультрафильтра $\frak F$ на полурешетке $\frak S_{\frak A}$ зафиксируем


некоторый ультрафильтр $\frak F'$


на $\omega$ такой, что $\frak F'\supseteq\frak F$. Так как алгебра


$\frak A$


богата, то существует дизъюнктное


вложение $\psi$ двоичного дерева $\frak D_2$ в полурешетку


$\frak S_{\frak A}=\la S^{\frak A},\cap\ra$. В


частности, существует


континуальная совокупность $A$ ультрафильтров на полурешетке


$\frak S_{\frak A}$ таких, что


для различных $\frak F$, $\frak G$ из $A$, для некоторых систем


$\frak T^{\frak F}_{\{\frak F,\frak G\}}$,


$\frak T^{\frak G}_{\{\frak F,\frak G\}}$


уравнений над $\frak A$


$$


S_{\frak A}(\frak T^{\frak F}_{\{\frak F,\frak G\}})\in\frak F,\quad


S_{\frak A}(\frak T^{\frak G}_{\{\frak F,\frak G\}})\in\frak G


\quad\text{и}\quad


S_{\frak A}(\frak T^{\frak F}_{\{\frak F,\frak G\}})\bigcap


S_{\frak A}(\frak T^{\frak G}_{\{\frak F,\frak G\}})=\varnothing.


$$


Через $A'$ обозначим совокупность $\{\frak F'\mid\frak F\in A\}$.





Для ультрафильтров $\frak F'$ на $\omega$ через


$\frak A'/\frak F'$ обозначим подалгебру


ультрастепени


$\frak A^{\omega}/\frak F'$, состоящую из элементов


$f/\frak F'$, где $f\in\frak A'$. Тогда


очевидно, что если


$$


\frak T^{\frak F}_{\{\frak F,\frak G\}}=


\begin{cases}


t^{1,\frak F}_{1,\{\frak F,\frak G\}}(x,a_0,a_1)=


t^{2,\frak F}_{1,\{\frak F,\frak G\}}(x,a_0,a_1),\\


\dots \\


t^{1, \frak F}_{n,\{\frak F,\frak G\}}(x,a_0,a_1)=


t^{2, \frak F}_{n,\{\frak F,\frak G\}}(x,a_0,a_1)


\end{cases}


$$


и


$$


\varphi^{\frak F}_{\{\frak F,\frak G\}}(x,y_0,y_1)=


\overset{n}{\underset{k=1}{\&}}{}


t^{1,\frak F}_{k,\{\frak F,\frak G\}}(x,y_0,y_1)=


t^{2,\frak F}_{k,\{\frak F,\frak G\}}(x,y_0,y_1),


$$


то


$$


\frak A'/\frak F'\vDash \varphi^{\frak F}_{\{\frak F,\frak G\}}


(g/\frak F',f_0/\frak F',f_1/\frak F')


$$


и


$$


\frak A'/\frak F'\vDash\lnot\,


\varphi^{\frak G}_{\{\frak F,\frak G\}}


(g/\frak F',f_0/\frak F',f_1/\frak F')


$$


для любых различных $\frak F$, $\frak G$ из $A$.





Далее, в силу универсальной аксиоматизируемости совокупности


$\frak K$-простых


алгебр алгебры $\frak A'/\frak F'$ просты и конечно-порождены. Покажем,


что существует


континуум ультрафильтров вида


$\frak F'\in A'$ таких, что алгебры $\frak A/\frak F'$


попарно невложимы


друг в друга, что и составит доказательство утверждения теоремы 2.





На множестве $A'$ введем отношение эквивалентности следующим образом:


$\frak F'\sim\frak G'$


тогда и только тогда, когда


$\frak A'/\frak F'\cong\frak A'/\frak G'$. Для каждого


$\frak F'\in A'$ множество


$\{\frak G'\in A'\mid\frak G'\sim\frak F'\}$ не


более чем счетно.





Действительно, допустим противное, и пусть для некоторого


$\frak F'\in A'$ множество


$\{\frak G'\in A'\mid\frak G'\sim\frak F'\}$


не счетно. В силу того, что все алгебры $\frak A'/\frak G'$ порождены


тремя элементами


$f_0/\frak G'$, $f_1/\frak G'$ и $g/\frak G'$, а сами эти алгебры счетны,


найдутся пара ультрафильтров $\frak G'_1\ne\frak G'_2$


и некоторый изоморфизм $\eta$ алгебр


$\frak A'/\frak G'_1$ и $\frak A'/\frak G'_2$ такие, что


$\eta(f_0/\frak G'_1)=f_0/\frak G'_2$,


$\eta(f_1/\frak G'_1)=f_1/\frak G'_2$,


$\eta(g/\frak G'_1)=g/\frak G'_2$. Но, с


другой стороны, эти ультрафильтры


$\frak G'_1$, $\frak G'_2$ отделимы некоторыми элементами


полурешетки $\frak S_{\frak A}$, т.е., как указано выше, найдется


система уравнений


$\frak T^{\frak G_1}_{\{\frak G_1,\frak G_2\}}$


над


алгеброй $\frak A$ такая, что для соответствующих построенных выше формул


$\varphi^{\frak G_1}_{\{\frak G_1,\frak G_2\}}$ имеет место


$$


\frak A/\frak G'_1\vDash


\varphi^{\frak G_1}_{\{\frak G_1,\frak G_2\}}


(g/\frak G'_1,f_0/\frak G'_1,f_1/\frak G'_1),


$$


но


$$


\frak A'/\frak G'_2\vDash\lnot\,


\varphi^{\frak G_1}_{\{\frak G_1,\frak G_2\}}


(g/\frak G'_2,f_0/\frak G'_2,f_1/\frak G'_2).


$$


Это противоречит существованию изоморфизма $\eta$. Таким образом,


действительно для любого $\frak F'\in A'\;$


$\bigl|\{\frak G'\in A'\mid\frak G'\sim\frak F'\}\bigr|\le


\aleph_0$, и, значит,


$|A'/\sim|=2^{\aleph_0}$.





Через $B$ обозначим $\{\frak G'\in A'\mid\text{ алгебра }


\frak A'/\frak G'\text{ не является изоморфно вложимой в алгебру }


\frak A\}$. Аналогично вычислению мощности множества $A'/\sim$


замечаем, что $|B/\sim|=2^{\aleph_0}$.


Более того, для любого кортежа $t\in\frak D_2$ имеет место равенство


$|\psi^*(t)\cap B/\sim|=2^{\aleph_0}$,


где $\psi$ ---


указанное в начале доказательства дизъюнктное вложение дерева


$\frak D_2$ в $\frak S_{\frak A}$, а


$\psi^*(t)=\{\frak G\in\frak S^*_{\frak A}\mid\psi(t)\in\frak G\}$. При


этом, очевидно, можно дополнительно считать выполненным следующее


условие: если $f$ --- любая бесконечная последовательность нулей и


единиц


и для $n\in\omega\;$ $f^n$ --- начальный интервал длины $n$


последовательности $f$, то


$\Bigl|\underset{n\in\omega}{\bigcap}{}\psi^*(f^n)\Bigr|=1$ и


$\underset{n\in\omega}{\bigcap}{}\psi^*(f^n)=\{\frak G\}$, где


$\frak G'\in B$.





Через $\frak D^n_2$ обозначим подмножество элементов дерева


$\frak D_2$, состоящее


из кортежей длины, не превышающей $n$. Через


$\frak D^*_2$ обозначим совокупность


последовательностей нулей и единиц длины $\omega$. Для


$f\in\frak D^*_2$, как и выше,


через


$f^n$ будем обозначать начальный интервал длины $n$ последовательности


$f$.





Индукцией по $n$ будем строить изоморфные вложения $\lambda_n$ частично


упорядоченных множеств $\frak D^n_2$ в дерево $\frak D_2$ таким образом,


что для $n<m\;$ $\lambda_m$


продолжает $\lambda_n$ и если для


$f\in\frak D^*_2\;$ $\lambda(f)=\frak G'_f$, где


$\frak G_f\in\underset{n\in\omega}{\bigcap}{}\psi^*(\lambda_n(f^n))$


(напомним, что $\psi$ --- дизъюнктное


вложение $\frak D_2$ в $\frak S_{\frak A}$), то для различных


$f_1,f_2\in\frak D^*_2$ алгебры


$\frak A'/\frak G'_{f_1}$ и $\frak A'/\frak G'_{f_2}$ невложимы друг в


друга.





Пусть


$$


\la p^0_0(y,z,x)=x, \ p^1_0(y,z,x)=y, \ p^2_0(y,z,x)=z\ra


,\dotsc,


\la p^0_n(y,z,x),p^1_n(y,z,x),p^2_n(y,z,x)\ra,\dotsc


$$


--- перечисление всех возможных троек термов от трех переменных сигнатуры


класса $\frak K$. Указанные выше вложения


$\lambda_n\;$ $\frak D^n_2$ в $\frak D_2$ будем строить так:


%%%%%%%%%%%% add \medskip





{\hfill


{\parbox{12cm}


{\it


для различных $t_1,t_2\in\frak D^n_2$ для любых


$\frak F_1$, $\frak F_2$ --- ультрафильтров


на $\frak G_{\frak A}$ таких, что


$\frak F_1\in\psi^*(\lambda_n(t_1))$,


$\frak F_2\in\psi^*(\lambda_n(t_2))$, отображения


$f_0/\frak F'_1\to p^0_m(f_0/\frak F'_2,f_1/\frak F'_2,g/\frak F'_2),$


$f_1/\frak F'_1\to p^1_m(f_0/\frak F'_2,f_1/\frak F'_2,g/\frak F'_2)$,


$g/\frak F'_1\to p^2_m(f_0/\frak F'_2,f_1/\frak F'_2,g/\frak F'_2)$


при $m\les n$ были бы непродолжимы до изоморфного вложения алгебры


$\frak A'/\frak F'_1$ в алгебру $\frak A'/\frak F'_2$.


}}


\hfill $(\ast)$}


%%%%%%%% add \medskip





Этого и будет достаточно, чтобы для различных


$f_1,f_2\in\frak D^*_2$ алгебры


$\frak A'/\frak G'_{f_1}$ и $\frak A'/\frak G'_{f_2}$ были бы


невложимы друг в друга.





Для $n=0$ положим $\lambda_1(\la 0\ra)=\la 0\ra$,


$\lambda_1(\la 1\ra)=\la 1\ra$. Очевидно, если


$\psi(\la i \ra)=\frak T_i$,


$$


\frak T_i=


\begin{cases}


t^1_1(x,a_0,a_1)=t^2_1(x,a_0,a_1),\\


\dots \\


t^1_n(x,a_0,a_1)=t^2_n(x,a_0,a_1)


\end{cases}


$$


и


$$


\varphi_{\frak T_i}(x,y_,y_1)=


\overset{n}{\underset{i=1}{\&}}{}


t^1_k(x,y_0,y_1)=t^2_k(x,y_0,y_1),


$$


то для любых $\frak G_1\in\psi^*(\la 0\ra)$,


$\frak G_2\in\psi^*(\la 1\ra)$


$$


\frak A'/\frak G'_1\vDash \varphi_{\frak T_0}


(g/\frak G'_1,f_0/\frak F'_1,f_1/\frak G'_1)


$$


и


$$


\frak A'/\frak G'_2\vDash\lnot\varphi_{\frak T_0}


(g/\frak G'_2,f_0/\frak G'_2,f_1/\frak G'_2),


$$


а значит, действительно, отображение $\eta$:


\begin{gather*}


f_0/\frak G'_1\to p^0_0


(f_0/\frak G'_2,f_1/\frak G'_2,g/\frak G'_2)=f_0/\frak G'_2,\\


f_1/\frak G'_1\to p^1_0


(f_0/\frak G'_2,f_1/\frak G'_2,g/\frak G'_2)=f_1/\frak G'_2,\\


g/\frak G'_1\to p^2_0


(f_0/\frak G'_2,f_1/\frak G'_2,g/\frak G'_2)=g/\frak G'_2


\end{gather*}


и обратное к нему непродолжимы до изоморфных вложений алгебры


$\frak A'/\frak G'_1$ в $\frak A'/\frak G'_2$ и


алгебры $\frak A'/\frak G'_2$ в $\frak A'/\frak G'_1$.





Пусть теперь уже построены отображения


$\lambda_1,\dotsc,\lambda_n$, удовлетворяющие указанным выше


индуктивным условиям, и пусть


$t_1,\dotsc,t_{2^n}$ --- все кортежи длины $n$, состоящие из


нулей и единиц. Фиксируем некоторое $j$ такое, что


$1<j\le 2^n$ и рассмотрим два


случая


\begin{itemize}


\item[а)] $\pi_{\psi(\lambda_n(t_j))}(p^0_{n+1}(f_0,f_1,g))


\ne\pi_{\psi(\lambda_n(t_j))}(p^1_{n+1}(f_0,f_1,g))$,\\


\item[б)] $\pi_{\psi(\lambda_n(t_j))}(p^0_{n+1}(f_0,f_1,g))


=\pi_{\psi(\lambda_n(t_j))}(p^1_{n+1}(f_0,f_1,g))$.


\end{itemize}





В первом случае для $i=0,1$ пусть $r^i_j$ является наименьшим


натуральным числом $r$ таким, что если


\begin{gather*}


\psi(\lambda_n(t_j){}\sphat\,


\underset{r\text{ раз}}{\la i,\dotsc,i\ra}{})=S_{\frak A}


(\frak T_i(x,a_0,a_1)),\\


\frak T_i(x,a_0,a_1)=


\begin{cases}


t^1_{1,i}(x,a_0,a_1)=t^2_{1,i}(x,a_0,a_1),\\


\dots \\


t^1_{n,i}(x,a_0,a_1)=t^2_{n,i}(x,a_0,a_1),


\end{cases}


\end{gather*}


то


\begin{multline*}


\psi(\lambda_n(t_j))\bigcap\bigcup\limits^n_{k=1}\bsl


t^1_{k,i}(p^0_{n+1}(f_0,f_1,g),p^1_{n+1}(f_0,f_1,g),


p^2_{n+1}(f_0,f_1,g))\ne\\


\ne t^2_{k,i}(p^0_{n+1}(f_0,f_1,g),p^1_{n+1}(f_0,f_1,g),


p^2_{n+1}(f_0,f_1,g))\bsr \ne\varnothing\tag{$\ast\ast$}


\end{multline*}


Прежде всего заметим, что такое $r$ найдется. В противном случае для


одного из $i$, равных нулю либо единице, относительно


$r\in\omega$ и любых систем


уравнений $\frak T^r_i(x,a_0,a_1)$ таких, что


$$


\psi(\lambda_n(t_1){}\sphat\,


\underset{r\text{ раз}}{\la i,\dotsc,i\ra}{}


=S_{\frak A}(\frak T^r_i(x,a_0,a_1)),


$$


имеет место включение


\begin{multline*}


\psi(\lambda_n(t_j))\subseteq\bigcap\limits^{n_r}_{k=1}\bsl


t^{1,r}_{k,i}(p^0_{n+1}(f_0,f_1,g),p^1_{n+1}(f_0,f_1,g),


p^2_{n+1}(f_0,f_1,g))=\\


= t^{2,r}_{k,i}(p^0_{n+1}(f_0,f_1,g),p^1_{n+1}(f_0,f_1,g),


p^2_{n+1}(f_0,f_1,g))\bsr,


\end{multline*}


где


$$


\frak T^r_i=


\begin{cases}


t^{1,r}_{1,i}(x,a_0,a_1)=t^{2,r}_{1,i}(x,a_0,a_1),\\


\dots \\


t^{1,r}_{n_r,i}(x,a_0,a_1)=t^{2,r}_{n_r,i}(x,a_0,a_1).


\end{cases}


$$


Но тогда если


$\{\frak F\}=\bigcap\limits_{r\in\omega}\psi^*(\lambda_n(t_1){}\sphat\,


\underset{r\text{ раз}}{\la i,\dotsc,i\ra}{})$,


то $\frak F'\in B$ и


отображение


$\eta:f_0/\frak F'\to p^0_{n+1}(f_0,f_1,g)\upharpoonright\psi(\lambda_n


(t_j))$,


$f_1/\frak F'\to p^1_{n+1}(f_0,f_1,g)\upharpoonright\psi(\lambda_n


(t_j))$,


$g/\frak F'\to p^2_{n+1}(f_0,f_1,g)\upharpoonright\psi(\lambda_n


(t_j))$


продолжимо


до изоморфного вложения алгебры $\frak A'/\frak F'$ в декартову степень


$\frak A^{\psi(\lambda_n(t_j))}$


алгебры $\frak A$.





Действительно, из равенства


$\Bigl|\bigcap\limits_{r\in\omega}\psi^*(\lambda_n(t_1){}\sphat\,


\underset{r\text{ раз}}{\la i,\dotsc,i\ra}{})\Bigr|=1$ очевидно


следует, что позитивная диаграмма


алгебры $\frak A'/\frak F'$ аксиоматизируема совокупностью равенств


$\bigcup\limits_{r\in\omega}\frak T^r_i(g/\frak F',f_0/\frak F',


f_1/\frak F')$, а


значит, и на


алгебре $\frak A^{\psi(\lambda_n(t_j))}$ истинна позитивная диаграмма


алгебры $\frak A'/\frak F'$ при


интерпретации порождающих алгебру $\frak A'/\frak F'$ элементов


$f_0/\frak F'$, $f_1/\frak F'$ и


$g/\frak F'$ элементами


$\pi_{\psi(\lambda_n(t_j))}(p^0_{n+1}(f_0,f_1,g))$,


$\pi_{\psi(\lambda_n(t_j))}(p^1_{n+1}(f_0,f_1,g))$,


$\pi_{\psi(\lambda_n(t_j))}(p^2_{n+1}(f_0,f_1,g))$


алгебры $\frak A^{\psi(\lambda_n(t_j))}$ соответственно.





С другой стороны, алгебра $\frak A'/\frak F'\;$ $\frak K$-проста, а


значит, поскольку


$\pi_{\psi(\lambda_n(t_j))}(p^0_{n+1}(f_0,f_1,g))\ne


\pi_{\psi(\lambda_n(t_j))}(p^1_{n+1}(f_0,f_1,g))$,


то отображение $\eta$ действительно продолжимо до изоморфного вложения


$\frak A'/\frak F'$ в


$\frak A^{\psi(\lambda_n(t_j))}$. Та же $\frak K$-простота алгебры


$\frak A'/\frak F'$


влечет при этом вложимость $\frak A'/\frak F'$ в алгебру


$\frak A$ в противоречие с тем, что $\frak F'$ является элементом


множества $B$. Таким


образом, указанные в условии $(\ast)$ числа $r^j_i$ найдутся. Пусть


$r=\max\limits_{1<j\les 2^n,\:i=0,1}r^j_i$.


Заметим теперь, что найдется кортеж $t'_j\in\frak D_2$, продолжающий кортеж


$\lambda_n(t_j)$ и


такой, что


\begin{multline*}


\psi(t'_j)\subseteq\bigcap\limits_{i=0,1}


\bigcup\limits^n_{k=1}\bsl


t^{1}_{k,i}(p^0_{n+1}(f_0,f_1,g),p^1_{n+1}(f_0,f_1,g),


p^2_{n+1}(f_0,f_1,g))\ne \\


\ne t^{2}_{k,i}(p^0_{n+1}(f_0,f_1,g),p^1_{n+1}(f_0,f_1,g),


p^2_{n+1}(f_0,f_1,g))\bsr,


\end{multline*}


где


$$


\psi(\lambda_n(t_1){}\sphat\,\underset{r\text{ раз}}{\la i,\dotsc,i\ra}{})


=S_{\frak A}(\frak T_i(x,a_0,a_1))


$$


и


$$


\frak T_i(x,a_0,a_1)=


\begin{cases}


t^1_{1,i}(x,a_0,a_1)=t^2_{1,i}(x,a_0,a_1),\\


\dots \\


t^1_{n,1}(x,a_0,a_1)=t^2_{n,i}(x,a_0,a_1).


\end{cases}


$$


Действительно, пусть никакого подобного кортежа $t'_j$ не нашлось и пусть


$f\in\frak D^*_2$ и $f$ продолжает кортеж $\lambda_n(t_j)$. Тогда для


любого $m\in\omega$ имели бы $\psi(f^m)\cap A\ne\varnothing$ и


$\psi(f^m)\cap\lnot A\ne\varnothing$, где


\begin{multline*}


A=\bigcap\limits_{i=0,1}\bigcup\limits^n_{k=1}\bsl


t^{1}_{k,i}(p^0_{n+1}(f_0,f_1,g),p^1_{n+1}(f_0,f_1,g),


p^2_{n+1}(f_0,f_1,g))\ne \\


\ne t^{2}_{k,i}(p^0_{n+1}(f_0,f_1,g),p^1_{n+1}(f_0,f_1,g),


p^2_{n+1}(f_0,f_1,g))\bsr.


\end{multline*}





В силу замкнутости $A$, $\lnot A$ и множеств $\psi(f_m)$ в компактном


пространстве $\St(\omega)$ имеют место и неравенства


$$


\bigcap_{m\in\omega}\psi(f^m)\bigcap A\ne\varnothing,\qquad


\bigcap_{m\in\omega}\psi(f^m)\bigcap\lnot A\ne\varnothing.


$$


В то же время


$\Bigl|\bigcap\limits_{m\in\omega}\psi^*(f^m)\Bigr|=1$


и пусть


$\bigcap\limits_{m\in\omega}\psi^*(f^m)=\{\frak F\}$,


где $\frak F\in\frak S^*_{\frak A}$. Если


$\frak G_1,\frak G_2\in\St(\omega)$ таковы, что


$\frak G_\in\bigcap\limits_{m\in\omega}\psi(f^m)\bigcap A$,


$\frak G_2\in\bigcap\limits_{m\in\omega}\psi(f^m)\bigcap\lnot A$,


то $\frak F\subseteq\frak G_1$ и


$\frak F\subseteq\frak G_2$.





Как уже отмечалось выше, позитивные диаграммы алгебр


$\frak A'/\frak G'$ и $\frak A'/\frak G_2$ полностью


определяемы совокупностью равенств


$\bigcup\limits_{m\in\omega}\frak T_m(g/\frak G_1,f_0/\frak G_1,


f_1/\frak G_1)$


и


$\bigcup\limits_{m\in\omega}\frak T_m(g/\frak G_2,f_0/\frak G_2,


f_1/\frak G_2)$


соответственно, где $\frak T_m(x,a_0,a_1)$ ---


система уравнений на $\frak A$, соответствующая кортежу


$f_m\in\frak D_2$ при отображении


$\psi$,


т.е. позитивные диаграммы алгебр $\frak A'/\frak G_1$ и


$\frak A'/\frak G_2$ совпадают при


замене $f_0/\frak G_1$, $f_1/\frak G_1$, $g/\frak G_1$ на


$f_0/\frak G_2$, $f_1/\frak G_2$, $g/\frak G_2$. С


другой стороны, т.к. $\frak G_1\in A$ и $\frak G_2\in\lnot A$, то


$$


\frak A'/\frak G_1\vDash\phi


(f_0/\frak G_1, f_1/\frak G_1, g/\frak G_1) \quad\text{ и }\quad


\frak A'/\frak G_2\vDash\lnot\phi


(f_0/\frak G_2, f_1/\frak G_2, g/\frak G_2),


$$


где


\begin{multline*}


\phi(y_0,y_1,x)=\underset{i=1,2}{\&}{}\bigvee^n_{k=1}


t^1_{k,i}(p^0_{n+1}(y_),y_1,x),p^1_{n+1}(y_0,y_1,x),


p^2_{n+1}(y_0,y_1,x))\ne \\


\ne t^2_{k,i}(p^0_{n+1}(y_0,y_1,x),p^1_{n+1}(y_0,y_1,x),


p^2_{n+1}(y_0,y_1,x)).


\end{multline*}


Полученное противоречие и доказывает существование указанного выше


кортежа $t'_j$.





В случае б), отмеченном в начале построения отображения


$\lambda_{n+1}$, т.е. когда


для $j$ ($1<j\les 2^n$) имеют место равенства


$$


\pi_{\psi(\lambda_n(t_j))}(p^0_{n+1}(f_0,f_1,g))=


\pi_{\psi(\lambda_n(t_j))}(p^1_{n+1}(f_0,f_1,g)),


$$


выбираем $t'_j$, равным $t_j$. Построенные таким образом для любого


$j$, $1<j\les 2^n$,


кортежи $t'_j$ обладают тем свойством, что для любых


$\frak F\in\psi^*(t_1{}\sphat\,\underset{r\text{ раз}}{\la


i,\dotsc,i\ra}{})$


и $\frak G\in\psi^*(t'_j)$ отображения


$\eta_e$


($e\les n+1$) такие, что


$f_0/\frak F'\to p^0_e(f_0/\frak G',f_1/\frak G',g/\frak G')$,


$f_1/\frak F'\to p^1_e(f_0/\frak G',f_1/\frak G',g/\frak G')$,


$g/\frak F'\to p^2_e(f_0/\frak G',f_1/\frak G',g/\frak G')$,


непродолжимы до изоморфного вложения


алгебры $\frak A'/\frak F'$ в алгебру $\frak A'/\frak G'$.





Аналогичные построения позволяют добиться того же и для любых


ультрафильтров


$\frak F\in\psi^*(t_1{}\sphat\,\underset{r\text{ раз}}{\la


0,\dotsc,0\ra}{})$,


$\frak G\in\psi^*(t_1{}\sphat\,\underset{r\text{ раз}}{\la


1,\dotsc,1\ra}{})$.


Продолжая эти же построения и далее для всех $t'_j$


($1<j\les 2^n$) на месте $t_1$, в конечном счете получаем необходимые


вложения $\lambda_{n+1}$


множества $\frak D^{n+1}_2$ в $\frak D_2$, для которых выполнено условие


$(\ast)$. Таким образом, как


уже было отмечено выше, для различных


$f_1,f_2\in\frak D^*_2$, если


$\frak G_{f_i}\in\bigcap\limits_{n\in\omega}


\psi^*(\lambda_n(f^n_i))$, то алгебры


$\frak A'/\frak G'_{f_1}$ и $\frak A'/\frak G'_{f_2}$


невложимы друг в друга.


\end{pf}





Из утверждений теорем 1 и 2 вытекает





\begin{cornonum}


Если $\frak K$ --- квазимногообразие, содержащее конечно-порожденную


богатую


$\frak K$-простую алгебру с одноэлементной подалгеброй, обладающее


свойством


Фрезера-Хорна, и класс $\frak K$-простых алгебр


$\forall$-аксиоматизируем, то отношения эпиморфности и вложимости на


$\frak K\;$ $\aleph_0$-идентичны.


\end{cornonum}





Условия следствия выполнимы, к примеру, для любого дискриминаторного


многообразия, содержащего конечно-порожденную богатую простую алгебру.





Любое не локально конечное дискриминаторное многообразие содержит простую


бесконечную конечно-порожденную алгебру, однако, как показывает следующий


пример, условие существования простой (не богатой) бесконечной


конечно-порожденной алгебры для справедливости заключения теоремы 1 не


достаточно.





Пусть $Z$ --- совокупность всех целых, а $N$ --- натуральных чисел, при


этом операция ${}'$ определена на $Z$ и $N$ как $x+1$. Пусть


$d(x,y,z)$ означает


функцию


дискриминатора на любом рассматриваемом множестве. Через


$\frak A_0$ обозначим


алгебру $\la N;{}',d\ra$, а через $\frak A_1$ --- алгебру


$\la Z;{}',d\ra$, тогда $\frak M=\frak M(\frak A_0)$ ---


не локально конечное


дискриминаторное многообразие. По теореме Йонссона совокупность подпрямо


неразложимых (простых) $\frak M$-алгебр содержится в классе


$SP_u(\frak A_0)$, т.е. в


классе,


состоящем из подалгебр ультрастепеней алгебры $\frak A_0$. Тем самым,


любая


простая $\frak M$-алгебра имеет вид дизъюнктного объединения алгебр вида


$\la N;{}'\ra$, $\la Z;{}'\ra$ с последующим определением функции $d$


как дискриминатора на этом


объединении. Итак, в качестве инвариантов для счетных простых


$\frak M$-алгебр


$\frak A$ выступает пара $\la u_{\frak A},v_{\frak A}\ra$, где


$u_{\frak A},v_{\frak A}\in\omega\cup\{\aleph_0\}$ и $u_{\frak A}$ ---


число дизъюнктных слагаемых вида


$\la N;{}'\ra$, а $v_{\frak A}$ --- вида $\la Z;{}'\ra$. Очевидно также,


что для простых счетных $\frak M$-алгебр


$\frak C_1$, $\frak C_2$ имеет место вложимость $\frak C_1$ в


$\frak C_2$ тогда и только тогда, когда


$v_{\frak C_1}\les v_{\frak C_2}$ и


$u_{\frak C_1}\les u_{\frak C_2}+(v_{\frak C_2}-v_{\frak C_1})$.


Здесь $\aleph_0-n=\aleph_0-\aleph_0=\aleph_0$ для любого $n\in\omega$.


Без труда замечается, что любая совокупность


попарно невложимых друг в друга счетных простых алгебр конечна.





\begin{thebibliography}{9}





\bibitem{P1}


Пинус А.Г. {\em Об отношениях вложимости и эпиморфности на


конгруэнц-дистрибутивных многообразиях} // Алгебра и логика. -- 1985. --


Т.\,25. -- \No\,5. -- С.\,588--607.


\bibitem{P2}


Pinus A.G. {\em Boolean constructions in Universal algebra}. -- Kluwer


Academ. Publ.: Dordrecht--Boston--London, 1993. -- 350~p.


\bibitem{P3}


Пинус А.Г. {\em Булевы конструкции в универсальной алгебре} // УМН. --


1992. -- Т.\,47. -- \No\,4. -- С.\,145--180.


\bibitem{P4}


Пинус А.Г. {\em Конгруэнц-дистрибутивные многообразия алгебр} // Итоги


науки и техн. ВИНИТИ. Алгебра. Топология. Геометрия. -- 1988. --


Т.\,26. -- С.\,45--83.


\bibitem{P5}


Пинус А.Г. {\em О простых счетных скелетах эпиморфности


конгруэнц-дистрибутивных многообразий} // Изв. вузов. Математика. --


1987. -- \No\,11. -- С.67--70.


\bibitem{Gor}


Горбунов В.А. {\em Характеризация резидуально малых квазимногообразий} //


ДАН СССР. -- 1984. -- Т.\,29. -- \No\,2. -- С.\,204--207.


\bibitem{P7}


Пинус А.Г. {\em О скелетах вложимости не локально конечных


дискриминаторных многообразий с бедной алгеброй} // Алгебра и логика. --


1994. -- Т.\,31. -- \No\,1. -- С.\,90--103.





\end{thebibliography}





\vskip 2cm





\post{\it Новосибирский государственный}{\it Поступила}


\post{\it технический университет}{20.04.1995}





\end{document}


